% !PRINTED 343
\weberpage{343}
Wenn aber $(1, \theta)$ eine Basis von $\Omega$ ist und, um beide Fälle
von §\,90, (9) zu umfassen,
\weq{2}{\theta = \frac{-\Delta+\sqrt{\Delta}}{2}}
angenommen wird, so ist im Falle a)
\[
\theta \text{ nicht kongruent mit einer Zahl } r,
\]
im Falle b) oder c)
\[
\theta \equiv r \pmod{\mathfrak{p}},
\]
wenn man also $-x = 2r + \Delta$ setzt, so folgt:

1. \gesperrt{Die notwendige und hinreichende Bedingung für
das Eintreten eines der Fälle b), c) ist die, daß es
eine rationale ganze Zahl gibt, die mit $\Delta$ zugleich
gerade oder ungerade ist, für die}
\weq{3}{\frac{x+\sqrt{\Delta}}{2} \equiv 0 \pmod{\mathfrak{p}}.}
Ist $\mathfrak{p}'$ das mit $\mathfrak{p}$ konjugierte Ideal, so ist
\weq{4}{\frac{-x+\sqrt{\Delta}}{2} \equiv 0 \pmod{\mathfrak{p}'},}
und daraus folgt
\weq{5}{x^2 \equiv \Delta \pmod{4p}.}
Ist umgekehrt die Kongruenz (5) durch $x = r$ befriedigt, so ist
\[
\frac{r+\sqrt{\Delta}}{2}\cdot\frac{-r+\sqrt{\Delta}}{2} \equiv 0 \pmod{\mathfrak{p}},
\]
und folglich genügt entweder $x = r$ oder $x = -r$ der Bedingung (3).

Die beiden Ideale $\mathfrak{p}, \mathfrak{p}'$ werden dann miteinander identisch
sein, wenn die beiden Zahlen (3) und (4) und folglich auch ihre
Summe $\sqrt{\Delta}$ durch $\mathfrak{p}$ und mithin $\Delta$ durch $p$ teilbar ist. Dann ist
$x \equiv 0$ oder $x \equiv p \pmod{2p}$ die einzige Wurzel von (5).

Wir haben also das Resultat:

2. \gesperrt{Der Primfaktor $\mathfrak{p}$ der natürlichen Primzahl $p$ ist
vom zweiten Grad, wenn die Kongruenz (5) keine
Lösung hat, vom ersten Grad, wenn (5) eine Lösung
hat, und $p$ ist das Quadrat von $\mathfrak{p}$, wenn (5) nur
eine Lösung hat, wenn also $p$ in $\Delta$ aufgeht.}

Das letztere ist in Übereinstimmung mit dem allgemeinen
Satze Bd.\,II, §\,174.
